% ============================================================================
% CHAPITRE 3 — Cadrage du Projet
% ============================================================================

\chapter{Cadrage du Projet}

Après l'analyse externe (Chapitre 2), ce chapitre pose le \textbf{cadre formel du projet}. Le tableau de cadrage constitue le « contrat de projet » — précis, daté, réaliste.

\section{Tableau de cadrage complet (Charte Projet)}

\rowcolors{2}{lightgray}{white}
\begin{table}[H]
\centering
\caption{Tableau de cadrage — Charte Projet}
\label{tab:cadrage}
\begin{tabularx}{\textwidth}{|l|X|}
\hline
\rowcolor{emsiblue}\textcolor{white}{\textbf{Rubrique}} & \textcolor{white}{\textbf{Contenu}} \\
\hline
Titre & Application Web de Gestion des Factures \\
Date début & Octobre 2025 \\
Date fin prévue & Janvier 2026 \\
Durée & 4 mois (16 semaines) \\
Contexte & PME/freelances marocains souffrant d'une gestion manuelle inefficace (Ch.1) \\
Objectif général & Développer une application web complète de gestion des factures, viable techniquement et économiquement \\
\hline
\rowcolor{lightblue}\multicolumn{2}{|l|}{\textbf{Objectifs SMART}} \\
\hline
Spécifique & CRUD factures/clients/articles, PDF, dashboard KPIs, rôles USER/ADMIN \\
Mesurable & 100\% fonctionnalités CDC ; 0 erreur calcul ; PDF $<$ 1s \\
Atteignable & Stack maîtrisée (React, Firebase) ; 4 mois ; 2 développeurs \\
Réaliste & Validation terrain : 83\% d'intention d'utilisation \\
Temporel & Livraison : jan 2026 ; soutenance : fév 2026 \\
\hline
\rowcolor{lightblue}\multicolumn{2}{|l|}{\textbf{Périmètre}} \\
\hline
Inclus & CRUD Factures/Clients/Articles/Catégories ; Calculs auto ; PDF + QR + Signature ; Dashboard ; USER/ADMIN ; Export Excel ; Email démo ; Multi-devise/société (interface) \\
Exclus & App mobile native ; Email automatique serveur ; JWT production ; Archivage cloud ; Multi-devise complet (V2) \\
\hline
\rowcolor{lightblue}\multicolumn{2}{|l|}{\textbf{Parties prenantes}} \\
\hline
Équipe & Abdessamad ELKHAL, Yahya FALHAOUI \\
Encadrants & Pr. Souissi Abdelmoughi, Pr. Hind Ouzif, Pr. Asmahane Tahiri \\
Utilisateurs & Comptables, freelances, dirigeants PME \\
Jury & Commission PFA EMSI \\
\hline
\rowcolor{lightblue}\multicolumn{2}{|l|}{\textbf{Contraintes}} \\
\hline
Délai & 4 mois non négociables \\
Budget & 0 MAD (outils open source, Firebase plan gratuit) \\
Tech & Hébergement local ports 3000/3001 ; auth localStorage (démo) \\
\hline
Livrables & Code source documenté ; rapport PFA ; démo fonctionnelle ; diagrammes UML \\
Critères réussite & Toutes fonctionnalités OK ; rapport $\geq$ 80 pages ; note $\geq$ 14/20 \\
\hline
\end{tabularx}
\end{table}

\section{Choix et justification du cycle de vie}

\subsection{Comparaison des modèles}

\rowcolors{2}{lightgray}{white}
\begin{table}[H]
\centering
\caption{Comparaison des modèles de cycle de vie}
\label{tab:cycle_vie}
\begin{tabularx}{\textwidth}{|l|X|X|l|}
\hline
\rowcolor{emsiblue}\textcolor{white}{\textbf{Modèle}} & \textcolor{white}{\textbf{Avantages}} & \textcolor{white}{\textbf{Inconvénients}} & \textcolor{white}{\textbf{Adapté}} \\
\hline
Cascade & Simple, documenté & Rigide, pas de retour arrière & ✗ \\
Modèle V & Tests associés & Peu flexible & ✗ \\
\textbf{Agile/Scrum} & \textbf{Itératif, adaptatif, livraisons fréquentes} & Disponibilité équipe requise & \textbf{✓} \\
Hybride & Mixte & Complexe à piloter & ✗ \\
\hline
\end{tabularx}
\end{table}

\subsection{Justification du choix Scrum}

\begin{enumerate}
    \item \textbf{Durée courte} (4 mois) : sprints de 2 semaines, livrables fréquents vérifiables.
    \item \textbf{Besoins évolutifs} : animations 3D et signature numérique ont émergé en cours de projet.
    \item \textbf{Binôme flexible} : user stories et backlog adaptés à 2 développeurs.
    \item \textbf{Validation continue} : chaque sprint livre un incrément démontrable.
\end{enumerate}

\section{Étude critique de l'existant}

\subsection{Solutions et limitations}

\rowcolors{2}{lightgray}{white}
\begin{table}[H]
\centering
\caption{Étude critique de l'existant}
\label{tab:existant}
\begin{tabularx}{\textwidth}{|l|X|X|}
\hline
\rowcolor{emsiblue}\textcolor{white}{\textbf{Solution}} & \textcolor{white}{\textbf{Description}} & \textcolor{white}{\textbf{Limitations}} \\
\hline
Excel/Word & Modèles manuels & Pas d'automatisation, erreurs, pas de PDF auto \\
QuickBooks & ERP USA & Anglais, 30-80\$/mois, trop complexe \\
Zoho Invoice & Cloud basique & Design daté, pas de rôles \\
Indy & Micro-compta freelance & Mono-utilisateur, pas de dashboard \\
Sage Maroc & ERP comptabilité & Prix élevé, installation locale \\
\hline
\end{tabularx}
\end{table}

\subsection{Critique}

Deux catégories de limites : outils \textbf{trop simples} (Excel) sans automatisation, et outils \textbf{trop complexes/coûteux} (QuickBooks, Sage) inadaptés aux PME. Aucun ne combine : accessibilité, design moderne, TVA différenciée, rôles USER/ADMIN, export Excel et archivage.

\section{Justification et différenciation}

\subsection{Avantages différenciateurs}

\begin{enumerate}
    \item \textbf{Design premium} : glassmorphism, animations 3D Framer Motion, QR code, signature
    \item \textbf{Spécialisation métier} : TVA par catégorie, remises ligne + globale, validation admin
    \item \textbf{Architecture dual-database} : Firebase (temps réel) + JSON Server (paramétrage)
    \item \textbf{Prix accessible} : 0 MAD démo, plans compétitifs vs concurrents
\end{enumerate}

\subsection{Positionnement}

La solution se positionne comme une \textbf{application web cloud-native, mobile-first}, développée selon les meilleures pratiques React (hooks, contextes, services séparés), offrant une expérience comparable aux meilleures SaaS du marché à coût minimal.
